% One-time imported draft. Editable native source; no Markdown regeneration.
\section{Step 22 Write an equivalent
  solution}\label{001-console-calculation-s22-step-22-write-an-equivalent-solution}

The free-writing contract describes behavior. An implementation can
compute the same result with a different scratch register or sequence of
additions. Teaching observations can name a register for a guided
exercise without making that register part of the final program
interface.

\subsection{Write or
  inspect}\label{001-console-calculation-s22-write-or-inspect}

Copy the starter into a new folder and implement the same contract with
a different correct arithmetic sequence. Test both input-domain
endpoints as well as the public examples. Explain why the transformation
preserves 2 \ensuremath{\times} x + 1.

\subsection{Expected
  observation}\label{001-console-calculation-s22-expected-observation}

The public checks should accept the alternative. Instruction count and
source spelling may differ; output and termination must agree.

\subsection{Check your
  understanding}\label{001-console-calculation-s22-check-your-understanding}

Which property must a free-writing solution preserve?

\begin{itemize}
  \tightlist
  \item
        \textbf{A} The specified input/output and termination behavior.
  \item
        \textbf{B} The exact private labels of a reference solution.
  \item
        \textbf{C} The exact scratch register used in the book.
\end{itemize}
